\begin{homeworkProblem}
  Dado $U \subset \mathbb{R}^m$ aberto, seja $f : U \to \mathbb{R}^n$ diferenciável no ponto $a \in U$. Prove que se $\lim v_k = v$ em $\mathbb{R}^m$ e $\lim t_k = 0$ em $\mathbb{R}$, então
  \begin{align*}
    \lim_{k\to\infty} \frac{f(a+t_k v_k)-f(a)}{t_k} = f'(a)\cdot v.   
  \end{align*}
  \begin{solution}
    Como $f$ é diferenciável no ponto $a$, se fixamos $t_kv_k\in\mathbb{R}^{m}$ taes que $a+t_kv_k\in U$, note que sem perdida de generalidade, suponha $|v_k|=1$, ja que se definimos $\overline{t}_k=t_k|v_k|$, $\overline{t}_k\to 0$ porque $t_k\to 0$ e $|v_k|\to|v|=l<\infty$, então usando a diferenciabilidade do $f$ temos 
    \begin{align*}
      f(a+t_kv_k)-f(a)=f^{\prime}(a)\cdot t_kv_k + r(t_kv_k).
    \end{align*}
    Portanto
    \begin{align*}
      \lim_{k \to \infty}\frac{f(a+t_kv_k)-f(a)}{t_k}&=\lim_{k \to \infty}f^{\prime}(a)\cdot v_k+\frac{r(t_kv_k)}{t_k}\\
      &=\lim_{k \to \infty}f^{\prime}(a)\cdot v_k+\frac{r(t_kv_k)}{t_k|v_k|}\\
      &=f^{\prime}(a)\cdot v.
    \end{align*}
  \end{solution}
\end{homeworkProblem}
